% ---- pg-565 （书页 553） ----

这里的拟合优度是指相似性与最终距离之间的单调关系.

第二种偏离度的度量是由塔卡杨(Takane)等人(见文献[31])所引进的, 已成为一个更受欢迎的准则. 对给定维数 $q$, 将这个量度记为 S 应力, 其定义为:
\phantomsection\label{eq:12-25}
\begin{equation}
S \text{ 应力} = \left[ \frac{\sum\limits_{i<k}\sum (d_{ik}^2 - \hat{d}_{ik}^2)^2}{\sum\limits_{i<k}\sum d_{ik}^4} \right]^{1/2} \tag{12-25}
\end{equation}
也就是说, S 应力是将\hyperref[eq:12-23]{式(12-23)}中的 $d_{ik}$ 和 $\hat{d}_{ik}$ 用它们的平方代表后所得到的量度. S 应力的值介于 0 和 1 之间. 典型的情况是: 此值小于 0.1 意味着给定的点结构对被考察的对象来说是一个好的几何表示.

一旦项目处于 $q$ 维空间, 它们的 $q\times 1$ 维坐标向量就被作为多元观测值来对待. 为了便于直观展示, 可以将 $q$ 维散点用它们的主成分坐标来表示(见第 8 章).

我们已将应力写为几何表示的维数 $q$ 的函数. 对每一个 $q$, 可以找到使应力达到最小的点结构. 随着 $q$ 的增加, 最小应力将在运算误差的范围内逐渐下降, 且当 $q=N-1$ 时达到零. 从 $q=1$ 开始, 可将应力$(q)$对 $q$ 作图. 这些点随 $q$ 的增加而呈下降排列. 若找到一个 $q$, 上述下降趋势到这一点开始接近水平状态, 即形成一个“肘”形曲线, 这个 $q$ 便是“最佳”维数.

整个多维标度变换算法可归纳成以下几个步骤:

1. 对 $N$ 个项目, 得出各对不同项目之间的 $M=N(N-1)/2$ 个相似性(距离). 将这些相似性按\hyperref[eq:12-21]{式(12-21)}那样排序. (距离则按从大到小的方式排序.) 如果相似性(距离)无法算出, 则必须指出它们的等级顺序.

2. 用一个尝试性 $q$ 维点结构计算点间距离 $d_{ik}^{(q)}$ 和数 $\hat{d}_{ik}^{(q)}$, 使后者满足\hyperref[eq:12-22]{式(12-22)}, 从而使\hyperref[eq:12-23]{式(12-23)}中的应力或\hyperref[eq:12-25]{式(12-25)}中的 S 应力最小化. (在计算机多维标度变换程序中, 常用一些能产生单调“拟合”距离的回归方法来确定 $\hat{d}_{ik}^{(q)}$ 值.)

3. 对给定的 $\hat{d}_{ik}^{(q)}$, 移动前一组点, 以便得到一个新的改进的点结构. (对固定的 $q$, 改进的点结构可通过将一般的函数最小化程序用于应力来确定. 这里应力被视作 $N$ 个项目的 $N\times q$ 个坐标的函数.) 对新的点结构可以计算新的点间距离 $d_{ik}^{(q)}$ 并进而得到新的 $\hat{d}_{ik}^{(q)}$, 因而有更小的应力. 以上过程不断重复, 直至得出最优表示(最小应力).

4. 将最小应力对 $q$ 作图, 从对图形的考察中选择最优维数 $q^{*}$. \hfill (12-26)

我们已假定最初的相似性数据是对称的($s_{ik}=s_{ki}$), 即没有结(两个或多个相似性相等), 也没有缺失数据. 克鲁斯卡尔(见文献[19,20])提出了一些在出现有非对称性、有结和(或)有缺失数据的情况时处理问题的方法. 此外, 现在有一些多维标度变换的计算机程序不仅可以处理欧氏距离, 而且可以处理一般的闵可夫斯基距离[见\hyperref[eq:12-3]{式(12-3)}].

以下三个例子用来说明以距离作为初始相似(异)性量度的多维标度变换.

\noindent\textbf{例 12.14}\quad (美国各城市的多维标度变换)

\hyperref[tab:12.7]{表 12.7} 给出选定的各对美国城市之间的航线距离. 因为城市自然位于一个二维空间之中(地球曲面一个接近平面的部分), 所以毫不奇怪, $q=2$ 的多维标度变换将这些城市表示为平面上的点, 因为这些点的相对位置大体与它们在地图上的相对位置一致. 注意, 如果表中的距离依从大到小的顺序排列, 即按从最不相似到最相似的顺序排列, 则第一个位置是 $d_{\text{波士顿, 洛杉矶}}=3\,052$.
